\centerline{\bf \Huge Разнобой по ТЧ}

\begin{flushright}
        \scriptsize{Эпизод 12. Ценность чуда}
\end{flushright}

\problem{1} 
Числа $d$ и $e$ являются корнями уравнения $x^2-b x+c=0$, причём все числа $b, c, d, e$ натуральны. Известно, что $b c d e=5202$. Какое наибольшее значение может принимать число $c$ ?

\problem{2} 
У Матвея есть $N$ квадратных плиток размером $1 \times 1$. Он может выложить из всех плиток прямоугольник, у которого короткая сторона больше 1, а длинная сторона на 11 длиннее короткой стороны. Матвей попробовал выложить из всех плиток квадрат, но у него осталось 8 плиток. Чему может равняться число плиток $N$ ?

\problem{3} 
На доске написаны пять последовательных натуральных чисел. Вася утверждает, что их произведение в 120 раз больше некоторого шестизначного числа, имеющего вид $A B A B A B$. То есть, у него в разрядах десятков, тысяч и сотен тысяч стоит ненулевая цифра $A$, а в остальных разрядах стоит цифра $B$, не обязательно отличающаяся от $A$. Найдите все возможные наборы чисел на доске, при которых слова Васи являются правдой.

\problem{4} 
Сумма четырёх попарно различных простых чисел равна произведению двух из них. Какое наименьшее значение может принимать эта сумма?

\problem{5} 
Натуральное число $M$ представили в виде произведения простых сомножителей. Затем каждый из них увеличили на 1 , и произведение стало равно $N$. Оказалось, что $N$ делится на $M$. Докажите, что если теперь разложить $N$ на простые множители и каждый из них увеличить на 1 , то полученное произведение будет делиться на $N$.


\problem{6}
На доске написаны два натуральных числа, одно из которых получается из другого перестановкой цифр. Может ли их разность равняться 2025? (Запись натурального числа не может начинаться с нуля).

\problem{7}
На доске записано пятизначное число. Петя стёр его первую и последнюю цифры, и число уменьшилось ровно в 46 раз. Какое число было записано изначально?

\problem{8}
Будем обозначать $\overline{a b c}$ число, состоящее из трёх различных цифр $a, b$ и $c$ в указанном порядке. Какое наибольшее значение может принимать НОД $(\overline{a b c}, \overline{b c a})$ ?

\problem{9}
Найти все натуральные числа $n$ такие, что все цифры в десятичной записи числа $6^n+1$ одинаковы.

\newpage

\hspace{3mm}

\problem{10}
У 10 детей есть несколько мешков с конфетами. Дети начинают делить конфеты между собой. Каждый по очереди забирает из каждого мешка свою долю и уходит. Доля вычисляется так: делим текущее число конфет в каждом мешке на число оставшихся детей (включая себя), если нацело не поделилось - округляем до целого в меньшую сторону. Может ли всем достаться разное количество конфет, (a) если мешков всего 8; (b) если мешков всего 9 ?

\problem{11}
На доску записали несколько (больше одного) последовательных натуральных чисел. Могло ли так случиться, что и сумма всех четных выписанных чисел - квадрат натурального числа, и сумма всех нечетных выписанных чисел - квадрат натурального числа?